\documentclass[10pt,a4paper]{amsart}
\usepackage[margin=19mm]{geometry}
\usepackage{amsmath,amssymb,mathtools}
\usepackage[scr=rsfs,cal=euler]{mathalfa}
\usepackage{xfrac}
\usepackage[unicode,hidelinks,bookmarks=false]{hyperref}
\newcommand{\ie}{i.e.\ }
\newcommand{\cf}{cf.\ }
\newcommand{\resp}{respectively\ }
\newcommand{\Subsection}{\subsection}
\providecommand{\nobreakdash}{\nobreak\hbox{-}}
\setlength{\emergencystretch}{2em}
\multlinegap0pt
\usepackage{amssymb}
\usepackage[shortlabels]{enumitem}
\newtheorem{lemma}{Lemma}
\title{Localization of Bergman Kernels and the Cheng--Yau Conjecture on Real Analytic Pseudoconvex Domains}
\author{Chin-Yu Hsiao, Xiaojun Huang, Xiaoshan Li}
\date{Source: arXiv:2604.05542v2 (CC BY 4.0)}
\begin{document}
\maketitle

\noindent\emph{Source and licence.} Localization of Bergman Kernels and the Cheng--Yau Conjecture on Real Analytic Pseudoconvex Domains. Version arXiv:2604.05542v2 (CC BY 4.0), distributed by arXiv under the Creative Commons Attribution 4.0 International licence, http://creativecommons.org/licenses/by/4.0/, as recorded in the arXiv licence metadata for this exact version and shown on the abstract page https://arxiv.org/abs/2604.05542v2. Full licence notice: \texttt{licenses/arxiv-2604.05542-CC-BY-4.0.txt}.\par\smallskip

\noindent\textit{Editorial scope.} Counted unit: the article's lemma 2-3-lem1 (source0.tex lines 1052--1055). The article's complete original proof of it is delivered with it (source0.tex lines 1056--1126, one \texttt{proof} environment of 71 lines). No mathematics is written, repaired or re-derived: the statement and the proof are the article's own lines, in the article's own notation, and no retained line is substituted or re-flowed.  No further source-local context is needed: every cross-reference of every retained line resolves inside the two delivered spans.  Nothing was omitted from any retained span, so the package carries no omission record.  No retained line contains a citation, so the wrapper carries no bibliography and no bibliography entry.  No retained line contains a figure environment, an image-inclusion command or drawing code, so no image is delivered and the package declares no asset dependency.  Editorial formatting only, in addition: an AMS \texttt{amsart} preamble that restates the article's own declarations and macros used by the retained lines, namely the environments \texttt{lemma} on separate counters, together with the standard packages those lines need.  The retained environments print in this document as Lemma 1 in order of appearance, whereas the article prints the same items with the numbering of its own full document.  The complete native source of the article as distributed in the arXiv e-print for this exact version is bundled unchanged as \texttt{sources/197961/source0.tex}.
\begin{lemma}\label{2-3-lem1}
Let $S_0\subset B_R(0)$ be a smooth germ of a complex analytic  hypervariety $S$ in $B_{R''}(0)$.
    For an open dense subset of $\{v\in\mathbb C ^{n}: \|v\|<\varepsilon(R)\}$, we have that $\mathcal R_{R, v}(S_0)$ is not a germ of $S$.
\end{lemma}

\begin{proof}
Let \(p_0 \in S_0\). First, suppose that near \(p_0\), \(S_0\) is the graph of a holomorphic function \(w = h(z)\), or equivalently, that the vector \(\frac{\partial}{\partial w}\big|_{p_0}\) is transverse to \(T_{p_0}^{1,0}S_0\). Taking \(v = 0\), the set \(\mathcal R_{R,0}(S_0)\) near \(\mathcal R_{R, 0}(p_0)\) is described by
\[
w = -\overline{h(z)} + P(z,\overline{z}),
\]
where 
\[
\mathcal R_{R,0}(z,w) = \bigl(z,\; -\overline{h(z)} + P(z,\overline{z})\bigr), ~(z, w)\approx p_0 ~\text{and}~ (z, w)\in S_0.
\]
Since \(-\overline{h(z)} + P(z,\overline{z})\) is not holomorphic in \(z\), the image \(\mathcal R_{R,0}(S_0)\) is not a complex submanifold near \(\mathcal R_{R,0}(p_0)\). Consequently, \(\mathcal R_{R,0}(S_0)\) cannot coincide with \(S\) in any neighbourhood of \(\mathcal R_{R,0}(p_0)\).

Now, suppose that for every point \(p\) near \(p_0\) we have \(\frac{\partial}{\partial w}\big|_{p} \in T_p^{1,0}S_0\). 
We may assume that near \(p_0\) the set \(S_0\) is described by
\begin{equation*}
z_1 = h(z',w),\qquad z' = (z_2,\dots ,z_{n}).
\end{equation*}
Since \(\frac{\partial}{\partial w}\) is tangent to \(S_0\) near \(p_0\), we have
\begin{equation*}
\frac{\partial}{\partial w}\bigl(-z_1+h(z',w)\bigr)\equiv 0,
\end{equation*}
and therefore
\begin{equation*}
h(z',w)=h(z')\quad\text{is independent of } w.
\end{equation*}
Thus, near \(p_0\) the set \(S_0\) can be written as
\begin{equation*}
z_1 = h(z'),\quad z' = (z_2,\dots ,z_n).
\end{equation*}
Consequently,
\begin{equation*}
\mathcal R_{R,v}(S_0)=\{(\widehat z,\widehat w):\widehat z=z+v(\widehat w-w),\;\widehat w=-\overline w+P(z+v(\widehat w-w),\overline z),\;z_1=h(z')\}.
\end{equation*}

Suppose further that, as germs, we have \(\mathcal R_{R,v}(S_0)\subset S\). Then $\mathcal R_{R, v}(S_0)$ is a germ of complex analytic hypervariety and we assume that $\mathcal R_{R, v}(S_0)=S$ near $\mathcal R_{R, v}(p_0)$.  Recall that \(\mathcal R_{R,v}(S_0)\) is given by the system
\begin{equation*}
\begin{split}
&\widehat z = z+v(\widehat w-w),\\
&\widehat w = -\overline w + P\bigl(z+v(\widehat w-w),\overline z\bigr),\\
&z_1 = h(z').
\end{split}
\end{equation*}
Choosing \(v_2=\cdots =v_{n}=0\), we obtain that 
\(\mathcal R_{R,v}(S_0)\) is described by
\begin{equation*}
\begin{split}
&\widehat z_1 = h(\widehat z_2,\dots ,\widehat z_{n}) + v_1(\widehat w-w),\\
&\widehat w = -\overline w + P\bigl(z+v(\widehat w-w),\overline z\bigr).
\end{split}
\end{equation*}
Then
\begin{equation*}
    v_1(\widehat w-w)=\widehat z_1-h(\widehat z_2, \cdots, \widehat z_n),\quad (\widehat z_1, \cdots, \widehat z_n, \widehat w)\in S,
\end{equation*}
and $v_1(\widehat w-w)$ as a function of $(\widehat{z},\widehat{w})\in S$ is independent of the parameter $v_1$. 
Here, we regard $w$ as a real analytic function in $(\widehat{z},\widehat{w})\in S$ with parameter $v_1$. Since 
$$\overline w = -\widehat w + P\bigl(\widehat{z}, \overline{\widehat{z_1}-v_1(\widehat{w}-w)},
\overline{\widehat{z_2}},\cdots,\overline{\widehat{z_n}}), \quad (\widehat z_1, \cdots, \widehat z_n, \widehat w)\in S,
$$
$w$ as a function on $S$ with parameter $v_1$ is independent of $v_1$, too.
Hence, it follows that over $S$
\[
\widehat w - w \equiv 0,
\]
and consequently on $S_0$
\[
w+\overline w = P(z,\overline z).
\]
Thus, for \((z,w)\) near \(p_0\) with \((z,w)\in S_0\), we have \((z,w)\in  M\). This implies \(S_0 \subset  M\), which contradicts the assumption  that \(M\) is of finite type.

Therefore, we have proved: If \(S\) is a complex analytic variety of codimension one in \(\mathbb C^{n+1}\) and \(S\cap B_R(0)\neq \emptyset\), then there exists a vector \(v\in\mathbb C^{n}\) with \(\|v\|\ll 1\) such that the intersection \(\mathcal R_{R,v}(S\cap B_R(0))\cap S\) has real codimension at least two in  $S$.
\end{proof}

\end{document}
